\section{Counting Sort: Loops in Sequence and Atoms for Hypotheses}\label{sec:cs}

\subsection{The program}
\begin{lstlisting}
method main(a: array<int, int>, n: int, k: int) returns (b: array<int, int>)
  requires 0 <= n && 0 < k
  requires keys_below(a, n, k)              // forall i :: 0 <= i < n ==> 0 <= a[i] < k
  ensures sorted(b, 0, n)
  ensures forall x: int :: count(x, selects(b, butlast(counter(0, 1, natural(n)))))
    == count(x, selects(a, butlast(counter(0, 1, natural(n)))))
{
  var c: array<int, int>; var i: int; var v: int; var s: int; var t: int;
  v := 0;
  while v < k { c := c[v := 0]; v := v + 1; }
  i := 0;
  while i < n { c := c[a[i] := c[a[i]] + 1]; i := i + 1; }
  s := 0; v := 0;
  while v < k { t := c[v]; c := c[v := s]; s := s + t; v := v + 1; }
  i := 0;
  while i < n { b := b[c[a[i]] := a[i]]; c := c[a[i] := c[a[i]] + 1]; i := i + 1; }
}
\end{lstlisting}
Four loops in sequence: the histogram is zeroed, filled, replaced by its prefix sums (the start position
of each value), and the elements are scattered to their start positions, which are incremented.

\subsection{The trace equation system}
Each loop contributes a trace whose current value is the initial value of the next; the simplifier has
already applied the introduction rules of \texttt{teq.dfy} (stores, consts, sums, selects) when the goal
is printed ($V, V_2, K_2$ are counters of $k'$ steps with $\intf(k') = k$, $K, I, K_3, I_3$ counters of
$n'$ steps with $\intf(n') = n$):
\begin{align}
  c_1 &= \cur(\mathsf{stores}(c_0, \bl{V}, \bl{\mathsf{consts}(k', 0)})) \label{eq:cs1}\\
  c_2 &= \cur(\tr{c_1}{\bl{K}, \bl{I}}{c, k, i \mapsto c[a[k] := c[a[k]] + 1]}) \label{eq:cs2}\\
  c_3 &= \cur(\mathsf{stores}(c_2, \bl{V_2}, \bl{S})) \qquad S = \sums(0, \mathsf{selects}(c_2, \bl{K_2})) \label{eq:cs3}\\
  b' &= \cur(\mathsf{stores}(b, \mapc{\bl{K_3}, \bl{C}, \bl{I_3}}{k, c, i \mapsto c[a[k]]}, \mathsf{selects}(a, \bl{I_3}))) \label{eq:cs4}\\
  C &= \tr{c_3}{\bl{K_3}, \bl{I_3}}{c, k, i \mapsto c[a[k] := c[a[k]] + 1]} \label{eq:cs5}
\end{align}
The scatter \eqref{eq:cs4} is a $\mathsf{stores}$ whose keys are the positions read from the table $C$ as
it is incremented, and whose values are the elements of $a$.

\subsection{The lemma chain}
The histogram is characterized at each value (the generic lemma \lemma{current\_trace\_add\_at} of
\texttt{teq.dfy} is for a fixed key; here the key is $a[k]$):
\begin{align*}
  &\lemma{current\_trace\_histogram}:\\
  &\quad \cur(\tr{c_0}{\bl{\cnt(0, 1, n)}, hs}{c, k, h \mapsto c[a[k] := c[a[k]] + 1]})[v]\\
  &\qquad = c_0[v] + \mathsf{count}(v, \mathsf{selects}(a, \bl{\cnt(0, 1, n)}))
\end{align*}
It fires on the reads $c_2[v]$ inside $\mathsf{selects}(c_2, \bl{K_2})$ of \eqref{eq:cs3}, with
$c_1[v] = 0$ from \eqref{eq:cs1} by \lemma{current\_stores} of \texttt{teq.dfy}, so the prefix sums
become
\[ S = \sums(0, \mapc{\bl{K_2}, \bl{V_2}}{v, h \mapsto \mathsf{count}(v, \mathsf{selects}(a, \bl{\cnt(0, 1, n')}))}). \]
The table of start positions is then named by a lemma whose left-hand side is this whole term:
\begin{align*}
  &\lemma{rank\_table\_def}:\\
  &\quad \cur(\mathsf{stores}(c_0, \bl{\cnt(0, 1, kk)}, \bl{\sums(0, \mapc{\bl{\cnt(0, 1, kk)}, hs}{v, h \mapsto {}\\
  &\qquad\qquad\qquad \mathsf{count}(v, \mathsf{selects}(a, \bl{\cnt(0, 1, n)}))})}))
    = \mathsf{rank\_table}(c_0, a, n, kk)
\end{align*}
which rewrites $c_3$ to $\mathsf{rank\_table}(c_2, a, n', k')$, the number of elements of $a[0..n)$ below
each value $v < k$ (and $c_2$ beyond). Naming the table is essential: the scatter lemma below needs the
bound $kk$ of the keys, and a side condition $\mathsf{ranks}(c_0, a, n, kk)$ with $kk$ open would be
matched by no fact, while $\mathsf{rank\_table}(c_2, a, n, kk)$ in the pattern binds $kk$. Finally the
scatter:
\begin{align*}
  &\lemma{sorted\_scatter}:\\
  &\quad s = \intf(n)\ \wedge\ \mathsf{keys\_below}(a, s, kb)\ \wedge\ kb = \intf(kk)
    \ \wedge\ Iv = \cnt(0, 1, n)\ \wedge\ hs_3 = \bl{Iv}\\
  &\quad\Longrightarrow\ \sorted(\cur(\mathsf{stores}(b_0,
    \mapc{\bl{\cnt(0, 1, n)}, \bl{T}, hs_2}{k, c, h \mapsto c[a[k]]},\\
  &\qquad\qquad\qquad\qquad\qquad \mathsf{selects}(a, hs_3))), 0, s)\ \Longleftrightarrow\ \mathsf{true}\\
  &\qquad \text{where } T = \mathsf{trace}(\mathsf{rank\_table}(c_2, a, n, kk),\\
  &\qquad\qquad\qquad \mathsf{for}\ \bl{\cnt(0, 1, n)}, hs\ \mathsf{do}\ c, k, h \mapsto c[a[k] := c[a[k]] + 1])
\end{align*}
and \lemma{count\_scatter} with the conclusion
\[ \mathsf{count}(x, \mathsf{selects}(\cur(\ldots), \bl{\cnt(0, 1, n)})) = \mathsf{count}(x, \mathsf{selects}(a, \bl{\cnt(0, 1, n)})). \]
The left-hand side binds $n := n'$, $kk := k'$, $hs_3 := \bl{I_3}$ and the rest; the counter of the
values' $\mathsf{selects}$ is tied by the open trace variable $Iv$ (Section~\ref{sec:merge:lemmas}), since $\bl{I_3}$ and
$\bl{\cnt(0, 1, n')}$ are not identified by a closed equation. The precondition on the keys is the atom
$\mathsf{keys\_below}$ of the specification with an open bound $kb$, matched against the fact and then
identified with $\intf(kk)$: a quantified hypothesis $\forall i.\ 0 \le i < n \Rightarrow 0 \le a[i] < k$
is matched by no side condition, open or closed, so the specification uses the atom (as
\texttt{count\_occurrences.dfy} uses the vocabulary of \texttt{teq.dfy} in its postcondition).

Soundness of \lemma{sorted\_scatter}: with the start positions correct for all keys, element $a[k]$ is
stored at (number of elements below $a[k]$) + (number of equal elements before position $k$), which is
a bijection of $0, \ldots, n - 1$ that is monotone in the value, so the result is sorted on $0, \ldots, n - 1$
and has the elements of $a$. Both are asserted, not proved; the lemma is the algorithm's correctness
argument in one statement, over the exact shape of the two last loops.

\subsection{The idea}
Loops in sequence compose through the current values: each loop's trace has the previous result as its
initial value, and the acceleration lemmas of \texttt{teq.dfy} (stores read at a position, sums, consts)
push the characterization of one loop into the next without any lemma specific to counting sort, up to
the table of start positions. The algorithm-specific knowledge is concentrated in naming the table and
in the scatter lemma, whose left-hand side is the composition of the last two loops. What could not be
expressed is a quantified precondition as a side condition; the specification's vocabulary must provide
an atom for it.
